
LLM-generated code that passes dense differential test suites --- 764 tests per task --- routinely violates formal contracts on the same contract-relevant inputs: programs returning the ground-truth output on contract-violating inputs nonetheless fail reference postconditions in 40\% of evaluations. This paper introduces an oracle isolation design that separates contract oracle semantic strength from input generation strategy, using ContractEval's contract-violating test inputs (CVTs) as a controlled fixed input set to compare differential and contract oracles on identical inputs. Across 364 HumanEval+/MBPP+ tasks and 3 LLM families (Experiment A), the oracle-isolation gap is 0.40 (95\% CI [0.358, 0.445]; Wilcoxon p = 5.88e-38; n = 10,432 triples), with a contract-unique mass of 0.40 on CVT inputs --- programs whose output equals the ground-truth canonical yet whose behavior violates the reference contract. Adaptive property-based testing via icontract-hypothesis adds an additional mean 0.0999 absolute violations beyond static oracle inputs (95\% CI [0.065, 0.133]; Wilcoxon p = 2.64e-22), consistently across all 5 tested model families (Experiment B; n = 17,226 triples). Two secondary hypotheses yielded negative results: AST-based contract richness correlates only weakly with oracle-isolation gap (Spearman $\rho$ = 0.136, not $\rho \geq$ 0.30), and cross-model contract-satisfaction variation is negligible at n = 5 (task-controlled gap = 0.007; $\Delta R^2$ = 0.004). These findings demonstrate that contract oracles and differential test oracles are qualitatively different evaluation instruments, and argue for oracle-type diversity as a first-class dimension in LLM code evaluation.

